% Source: sources/returns/S6_04_MAPS_MANIFOLDS/sections/manifolds.tex:324–397
\paragraph{Theorem D.1.6 (Proved: the integral cohomology ring).}
Suppose $d>1$, so in this construction $d\geq5$. Number the cusp
components $D_j$, $j\in\mathbb Z/d$, by the positive $e_2$ deck
translation, and give all complex subvarieties their complex orientations.
Write $x_j=\operatorname{PD}[D_j]\in H^2(\widetilde X_p;\mathbb Z)$.
These classes generate, with the sole relation $\sum_jx_j=0$.
There is a distinguished generator $\kappa$ of the torsion subgroup of
$H^4(\widetilde X_p;\mathbb Z)$, namely the dual of the complex-oriented
ruling class $H-E_1$ in any normalized cusp component. The products are
\[
 x_i x_j=\begin{cases}
 -2\kappa,&i=j,\\
 \kappa,&i-j=1\text{ or }-1\pmod d,\\
 0,&\text{otherwise}.
 \end{cases}
\]
Together with the orientation pairing between the free parts of $H^2$
and $H^4$, these are every product in positive degree. For $d=1$ the
ring is the ring of $S^6$.
\begin{proof}
The relative cohomology group $H^2(N,M)=H_4(N)=\mathbb Z^d$ is
generated by the Thom classes of the irreducible Cartier divisors $D_j$.
This remains true for their self-identifications, since their fundamental
cycles generate $H_4(N)$ and Poincare--Lefschetz duality applies to the
smooth tube. The boundary homomorphism from $H^1(\widetilde X^\circ)$
takes its primitive generator to $(\pm1,\ldots,\pm1)$, with a common
sign. Indeed the toric meridians at adjacent height-one rays differ by
an angular $\widehat w$ or $\widehat\delta$ circle, and both have
zero image under $\Theta$; D.1.4 evaluates any one of them as $p'=\pm1$.
Consequently $\mathbb Z^d/\mathbb Z(1,\ldots,1)$ injects into
$H^2(\widetilde X_p)$. The cokernel embeds in the torsion-free group
$H^2(\widetilde X^\circ)=\mathbb Z^2$ and has rank zero by D.1.3.
It is therefore zero. This proves the assertion about the generators
integrally, including saturation.

For clarity we compute the relevant curve relations before their
images in the whole manifold. The first opposite pair on component
$j$ identifies $E_{1j}$ with $H_j-E_{2j}-E_{3j}$. The other two
opposite pairs identify $E_{2j}$ with $E_{2,j+1}$ and $E_{3j}$ with
$E_{3,j+1}$, after using the first relation. These are precisely the
transpose relations of $B(t)$ in D.1.2. Thus the ruling classes
$f_j=H_j-E_{1j}=E_{2j}+E_{3j}$ all have the same image $f$ in
$H_2(N)$. Components $j$ and $j+1$ intersect along the two external
boundary curves, with multiplicity one and total class $f$; no
nonadjacent components meet. Hence $x_jx_{j+1}=\operatorname{PD}(f)$.
Multiplying $\sum_ix_i=0$ by $x_j$ gives $x_j^2=-2\operatorname{PD}(f)$.

It remains to verify that $f$ is a torsion generator, rather than a
proper multiple. The torus $\widehat u\wedge\widehat\delta$ specializes
to one of the two external boundary-curve classes, up to sign, and
$\widehat u\wedge\widehat w$ specializes to their signed sum. One
can check this in a single strip of the covering tiling: the base
circle has direction $B_0\widehat u=e_1$; it crosses the two edges
with stabilizer directions $e_2-e_1$ and $-e_2$, and its degree on
each surviving angular circle is the determinant with that direction.
These determinants are units. This is the integral annulus-collapse
calculation of AC26 (7.9), applied to the unchanged $e_1$ period.
In the first finite piece, $\widehat u\wedge\widehat\delta$ has
coordinates $(0,0)$ in the dual bases of D.1.4, whereas
$\widehat u\wedge\widehat w$ has coordinates $(0,1)$.
The functional $\Phi=(4d,2,3d,2)$ therefore evaluates their images
as $0$ and $2$, respectively. It follows that one external curve
has image zero and the other has image $\pm2$ in
$\mathbb Z/d\subset H_2(\widetilde X_p)$. Their complex-oriented
sum $f$ is consequently a generator, since $d$ is odd. Its sign
defines the stated $\kappa$, so the intersection formulas themselves
have no remaining sign ambiguity.

All products involving a torsion class and landing in $H^6=\mathbb Z$
vanish. The groups in degrees one and five vanish, and a product of
two degree-three classes is torsion and hence zero. Finally Poincare
duality gives the perfect integral pairing of the free degree-two
and degree-four groups. This exhausts all degree possibilities.
\end{proof}
